\section{Fields on the round sphere}\label{sec:round}
First consider the geometric interpretation. Suppose that $F:S^n\to\R^d$
is tangent and equipped with a PR modulus $\nu$ for $\ell^1$ distance.
Suppose its values at rational unit vectors are given to any precision
by a fixed PR rational sampler. These are data one can supply, for
example, for expressions built from rational arithmetic, bounded
reciprocals and square roots on bounded intervals. The value at a
specified precision is understood to have $\ell^1$ error at most
$2^{-h}$. A coordinatewise evaluator can meet this convention by
requesting $\clog(d)$ additional bits.

For $p,q\in P_\Q$ let $\rho=\sqrt{p\cdot p}$,
$\eta=\sqrt{q\cdot q}$. Both belong to $[1/\sqrt d,1]$.
The reverse triangle inequality bounds $|\rho-\eta|$ by $\len{p-q}$.
Subtracting $p/\rho$ and $q/\eta$ and using $\len q=1$ shows
\[
 \len{\rad(p)-\rad(q)}
 \le \left(\frac1{\rho}+\frac1{\rho\eta}\right)\len{p-q}
 \le 2d\,\len{p-q}.
\]
This explains the precision shift between the two geometries. The
construction below gives a rational unit point $R(p,a)$ within
$2^{-a}$ of $\rad(p)$. Evaluate the round sampler to precision $h+2$
at $R(p,\nu(h+2))$ to obtain $A(p,h)$. Its distance from
$F(\rad(p))$ is at most $2^{-(h+1)}$. We may take
$\mu(h)=\nu(h)+\clog(2d)$. The two evaluation errors and the modulus
give (U), common-point errors give (C), and
$p\cdot F(\rad(p))=0$ gives (T). Thus Theorem~\ref{thm:main}
produces the asserted bound on the actual field value. The last
subsection gives a formulation of this transfer entirely in arithmetic.

\subsection{A rational witness on the round sphere}\label{sec:rational-round}

Suppose the field $F$ on the Euclidean unit sphere has an $\ell^1$
modulus $\nu$, so that input distance at most $2^{-\nu(j)}$ gives
output distance at most $2^{-j}$. Apply the direction-witness result at
precision $k+1$. It gives a rational $p$ of $\ell^1$ norm one whose
unit direction $x=p/\|p\|_2$ satisfies $\|F(x)\|_1\leq2^{-(k+1)}$.
We construct a rational unit vector $z$ with $\|F(z)\|_1\leq2^{-k}$.
The construction also bounds its denominator.

Choose the first $i$ maximizing $|p_i|$, put $q=|p_i|\geq 1/d$,
and let $\sigma$ be the sign of $p_i$. For a vector $u$
indexed by the $n=d-1$ remaining coordinates, write $U=\sum_{j\ne i}u_j^2$.
The inverse stereographic chart is
\[
  \Phi(u)_j=\frac{2u_j}{1+U}\quad(j\ne i),\qquad
  \Phi(u)_i=\sigma\frac{1-U}{1+U}.
\]
The identity $4U+(1-U)^2=(1+U)^2$ proves that the output lies on
the unit sphere. If $\rho=\sqrt{p\cdot p}$, the chart parameter of
$x$ is $u_j=p_j/(\rho+q)$. Indeed $U=(\rho-q)/(\rho+q)$, and
substitution gives $\Phi(u)=x$.

Put $a=\nu(k+1)$ and
\[
   t=a+\lceil\log_2(2d(d^2+d))\rceil,\qquad M=2^t,\qquad h=1/M .
\]
Here the integer ceiling logarithm is computed by bounded search.
Bisect $[0,1]$ exactly $t$ times, comparing each midpoint squared with
$p\cdot p$. Retain the upper half if the square is smaller, and the
lower half otherwise. The upper endpoint $b$ satisfies
$\rho\leq b\leq\rho+h$. Form $p_j/(b+q)$ and round each coordinate
toward zero to the nearest multiple $v_j=m_j/M$ of $h$ in that direction.
The total change from replacing $\rho$ by $b$ is at most
\[
  \sum_{j\ne i}\frac{|p_j|(b-\rho)}{(\rho+q)(b+q)}
      \leq d^2 h .
\]
Rounding adds at most $dh$ in $\ell^1$. Set $z=\Phi(v)$.
Expanding the chart formulas gives the chord identity
\[
 \|\Phi(v)-\Phi(u)\|_2^2=
 \frac{4\|v-u\|_2^2}
      {(1+\sum_{j\ne i}v_j^2)(1+\sum_{j\ne i}u_j^2)}.
\]
Thus $\|z-x\|_1\leq 2d\|v-u\|_1\leq2d(d^2+d)h\leq 2^{-a}$,
using $\|y\|_1\leq\sqrt d\,\|y\|_2$ and $\sqrt d\leq d$.
The modulus and the triangle inequality give the claimed bound on $F(z)$.
All computed quantities use rational operations. The occurrences of
$\rho$ and $u$ identify the direction whose distance the construction
certifies. The same bounds hold for their Cauchy names by the rational
bisection enclosures.

The common denominator $M^2+\sum_{j\ne i}m_j^2$ of the displayed
chart coordinates is at most $2M^2$, since $|v_j|\leq |u_j|$ and
$\sum u_j^2\leq1$. Consequently the possible outputs at precision $k$
belong to an explicit finite grid on the round sphere.


One also obtains a rational round witness directly from the arithmetic
name of Definition~\ref{def:name}, using only its modulus $\mu$.
For $q=W(k)$ and $t=k+6$, take $z=R(q,\mu(t)+1)$ and $w=z/\len z$.
Normalization $y\mapsto y/\len y$ is $2$-Lipschitz in $\ell^1$ on
the Euclidean unit sphere: subtract the two fractions, apply the
reverse triangle inequality, and use $\len y\ge1$. It follows that
$\len{w-q}\le2^{-\mu(t)}$, with both $w,q$ rational. By (U) and (C),
for $j\ge t$ the name of the field at the rational round point $z$ obeys
\[
       \len{A(w,j)}
        \le (3/4+3/64+2/64)\,2^{-k}
        =\tfrac{53}{64}\,2^{-k}.
\]

\subsection{A first-order statement using only rational unit points}
Here is a direct arithmetic proof for a field whose only input interface
is a PR sampler $B(z,h)$ on rational vectors satisfying $z\cdot z=1$.
Give it (C), (U), (T) on that domain, with $\nu$ in place of $\mu$ and
$z$ in the tangency test. There are no real quantifiers in this
hypothesis. We construct a PR rational $z=z(k)$ such that iPRA proves
\[
 z\cdot z=1,\qquad
 \forall j\ge k+10\quad \len{B(z,j)}\le 2^{-k}.                 \tag{R}
\]
In particular the actual real vector named by $B(z,-)$ has that bound.

Let $E(a)=2(2^{a+\clog(2d(d^2+d))})^2$, the denominator bound
for $R(p,a)$ just proved. Compute a uniform value bound once. Enumerate
all rational unit points written with a common denominator
$b\le E(\nu(0))$, $1\le b$, and numerators in $[-b,b]$.
Let $M$ be $5$ plus the ceiling of the largest $\ell^1$ norm of
$B(z,0)$ in this list. The axes ensure the list is nonempty.
Every rational unit $w$ has rational parameter $w/\len w$.
Its rounding at $a=\nu(0)$ is within $2^{-\nu(0)}$ and occurs on
the list. From (U) and (C) at scale zero we obtain
$\len{B(w,j)}\le M$ for every $j$. This derives the bound from the
input data.

For each $h$ put $s_h=h+4$ and choose the monotone exponent
\[
 a_h=2+\max\!\left(\max_{i\le s_h}\nu(i),\,
                      s_h+\clog(M+1)\right).
\]
Define $R_h(p)=R(p,a_h)$ and
\[
 A(p,h)=B(R_h(p),s_h),\qquad
 \mu(h)=\nu(s_h)+2+\clog(2d).
\]
If $j\ge h$, the round inputs $R_j(p)$ and $R_h(p)$ are at distance at
most $2\,2^{-a_h}\le2^{-\nu(s_h)}$. Compare the two samples first
by round (U) at $s_h$ and then by round (C) at $R_j(p)$. The result is
at most $5\,2^{-s_h}\le 2\,2^{-h}$. For (U), our $\mu(h)$ guarantees
that the exact directions are at distance at most $2^{-\nu(s_h)}/4$;
the two roundings add at most $2\,2^{-a_h}$. Round (U) then bounds the
sample difference by $3\,2^{-s_h}$. Finally write $x=\rad(p)$.
Since $p=\rho x$, $\rho\le1$, the tangency residual is bounded by
\[
 |p\cdot A(p,h)|
 \le |R_h(p)\cdot B(R_h(p),s_h)| + 2^{-a_h} M
 \le 2^{-s_h}+2^{-a_h} M\le2^{-h}.
\]
This verifies (T). Every estimate involving $\rho$ can equivalently be
read using rational upper and lower bisection endpoints. To eliminate
such an auxiliary quantity in an inequality between rationals, establish
the inequality with error $C2^{-b}$ at each precision $b$, with $C$
given by the displayed denominator bounds. If the rational inequality
fails, choose $b$ by bounded comparison with its positive rational
gap. This proves the exact rational inequality in iPRA. For example,
$\rho\eta\ge1/d$ follows directly by squaring, and dividing in a
bisection estimate may start with the lower bound $1/d$ for each
length.

Apply Theorem~\ref{thm:main} to this $A$ and take the resulting $q$.
The point $z=R_t(q)$, $t=k+6$, is rational and satisfies $z\cdot z=1$.
Its observed value is exactly $B(z,t+4)=A(q,t)$. At every $j\ge t+4$
the bound in (R) follows by adding $2\,2^{-(t+4)}$ to
$3\,2^{-k}/4$. Its common denominator is at most $E(a_t)$. This proof
of (R) operates on number codes, with the real sphere serving as the
interpretation of the rational computations.
